% Source: https://artofproblemsolving.com/community/c3h1254374_equal_angles_in_a_pentagon
% Stable identity: aops-topic-1254374-q01
The interior angles of a pentagon total $540^\circ$, so each of the other four is $\dfrac{540-100}{4}=110^\circ$.

\begin{asy}
usepackage("amsmath");
size(4.5cm);
defaultpen(fontsize(10pt));
pair Y=(0,0),X=(3,0);
pair W=X+3*dir(70); pair Z=Y+3*dir(110);
pair V=extension(W,W+dir(70+80),Z,Z+dir(110-70));
pair Q=extension(Y,Z,V,V+(X-Y));
draw(V--W--X--Y--Z--cycle,linewidth(1)); draw(Z--Q,dashed); draw(Q--V,blue);
label("$V$",V,N);label("$W$",W,E);label("$X$",X,SE);label("$Y$",Y,SW);label("$Z$",Z,W);label("$Q$",Q,NW);
label("\mbox{}",(min(currentpicture,user=true)+max(currentpicture,user=true))/2,black); // keeps the text after the figure black when LaTeX typesets the labels inline
\end{asy}

In triangle $QVZ$:
\begin{itemize}
\item $\angle QZV=180^\circ-\angle YZV=70^\circ$ (supplementary angles on line $YQ$);
\item since $QV\parallel YX$ and $V$, $X$ lie on the same side of line $YQ$, the co-interior angles $\angle VQY$ and $\angle QYX$ add to $180^\circ$, so $\angle VQZ=180^\circ-110^\circ=70^\circ$.
\end{itemize}
Hence $\angle QVZ=180^\circ-70^\circ-70^\circ=\boxed{40^\circ}$.
